%% P17 --- Hesjan, krzywizna i rozwinięcie Taylora
%%
%% Bliźniak programs/en/P17-hessian-taylor.tex. Ramka w ramkę, makro w makro,
%% liczba w liczbę; tools/parity.py porównuje oba pliki.
%%
%% UWAGA DLA TŁUMACZA: C4, C8 i C12 porównują znaczniki strukturalne, wzory i
%% literały liczbowe PO KOLEI. Zdanie polskie musi nieść te same wzory i te
%% same cyfry w tej samej kolejności, a nie tylko te same. Cyfra zostaje
%% cyfrą: \enquote{dwukrotnie} zamiast $2$ psuje C8 i C12 naraz. I odwołanie
%% do programu prowadzi zdanie, zamiast stać za wzorem.
%%
%% I skrót jest częścią strumienia znaczników: tam, gdzie angielski pisze
%% samo \enquote{P15} zamiast Program~\ref{}, polski też musi, bo C4 i C14
%% liczą odwołania. Wygląda to na wybór stylistyczny, a jest strukturalny.

\program[Hesjan i krzywizna]{Hesjan, krzywizna i rozwinięcie Taylora}
\label{prog:P17}
\index{hesjan}
\index{krzywizna}

\begin{outcomes}
\outcome{1--8}{powiedzieć, czym jest model Taylora drugiego rzędu i w jakim
  sensie jest najlepszym modelem lokalnym}
\outcome{9--15}{zapisać macierz drugich pochodnych i powiedzieć, dlaczego
  Program~\ref{prog:P10} stosuje się do niej bez zmian}
\outcome{16--23}{wyprowadzić największy rozmiar kroku, jaki toleruje funkcja
  kwadratowa, i powiedzieć, który kierunek go ustala}
\outcome{24--30}{powiedzieć, co dałoby użycie krzywizny wprost i ile by to
  kosztowało}
\outcome{31--38}{powiedzieć dokładnie, które twierdzenia o ostrych i płaskich
  minimach mówią o funkcji, a które o układzie współrzędnych}
\end{outcomes}

\begin{quiz}
\item Znasz wartość funkcji, nachylenie i krzywiznę w punkcie. Co najlepszego
  możesz powiedzieć o niej w pobliżu? \teachesat{1--4}
  \answerto{Parabola \dash{} model Taylora drugiego rzędu. Jej błąd maleje
  jak sześcian odległości, gdzie błąd stycznej maleje jak kwadrat, i to
  właśnie znaczy tu \enquote{najlepszy}: rząd, a nie rozmiar.}
\item Program~\ref{prog:P15} pokazał, że każdy kierunek jest mnożony przez
  $1 - \eta\lambda$ w każdym kroku. Jaki pułap kładzie to na $\eta$ i czy sam
  pułap jest użyteczny? \teachesat{16--17}
  \answerto{$\eta$ musi być ściśle poniżej $2/\lambda$, bo spacer pozostaje
  ograniczony dokładnie wtedy, gdy $\lvert 1 - \eta\lambda \rvert < 1$
  \dash{} nie ma więc największej użytecznej wartości, jest tylko granica,
  której zbiór nie zawiera. Dokładnie przy $2/\lambda$ czynnik wynosi $-1$ i
  współrzędna nigdy się nie kurczy, co rozpisuje ćwiczenie testowe 1. Przy
  wielu kierunkach wygrywa
  najmniejsze z tych ograniczeń, czyli $2/\lambda_{\max}$.}
\item Który kierunek doliny decyduje o rozmiarze kroku, a który musi z tym
  żyć? \teachesat{18--19}
  \answerto{Najbardziej stromy ustala pułap, bo to on rozbiega się pierwszy;
  najbardziej płaski musi z tym żyć, a jego własne ograniczenie jest większe
  o wskaźnik uwarunkowania.}
\item Metoda Newtona używa krzywizny, by wybrać krok osobno dla każdego
  kierunku. Dlaczego nie używa się jej do trenowania?
  \teachesat{26--27}
  \answerto{Bo hesjan ma jeden wpis na każdą \emph{parę} parametrów, a
  rozwiązanie układu z nim jest jeszcze gorsze. Przy realnym rozmiarze modelu
  sama macierz jest większa niż jakakolwiek istniejąca pamięć.}
\item Czy \enquote{to minimum jest ostrzejsze niż tamto} to fakt o
  wytrenowanym modelu? \teachesat{31--33}
  \answerto{Nie samo z siebie. Przeskalowanie parametru zmienia krzywiznę,
  zostawiając każde wyjście identyczne, więc goła liczba opisująca ostrość
  jest faktem o układzie współrzędnych. To, co przeżywa reparametryzację,
  trzeba wypowiedzieć ostrożnie, a sekcja 5 mówi, które części przeżywają.}
\end{quiz}

% ============================================================
\section{Najlepszy model lokalny}
\label{sec:P17-taylor}
\index{rozwinięcie Taylora}

\begin{fr}
Program~\ref{prog:F11} przybliżał krzywą w pobliżu punktu jej styczną, a
Program~\ref{prog:P15} zrobił to samo w wielu wymiarach za pomocą gradientu.
Styczna jest \emph{modelem}: zgadza się z funkcją w punkcie i w nachyleniu, a
wszędzie indziej się z nią rozmija.

Masz wartość, nachylenie, a teraz i tempo, w jakim zmienia się samo
nachylenie. Co zbudowałbyś z tych trzech rzeczy i jaki miałoby to kształt?
\dotline
\nextframe
\end{fr}

\begin{fr}
\begin{ansblock}
Parabolę \dash{} wartość plus nachylenie razy odległość plus połowa
krzywizny razy odległość do kwadratu.
\end{ansblock}

\[ f(x_0 + h) \;\approx\; f(x_0) + f'(x_0)\,h + \tfrac{1}{2}f''(x_0)\,h^{2} \]

To jest \textbf{model Taylora drugiego rzędu}, a $\tfrac{1}{2}$ nie jest
sztuczką: to ono sprawia, że druga pochodna modelu równa się drugiej pochodnej
funkcji, co jest całym powodem dokładania tego składnika.

\mermaidfig{p17-three-models}{Stały, liniowy, kwadratowy \dash{} jedna zgodność
więcej za każdym razem.}{trzy modele lokalne}
\end{fr}

\begin{fr}
Teraz teza, która czyni ten model wartym posiadania, a łatwo ją wypowiedzieć
niedbale.

Kusi, by powiedzieć, że parabola jest \enquote{dokładniejsza} niż styczna. To
prawda i to nie jest użyteczne stwierdzenie. Jakie jest użyteczne?
\nextframe
\end{fr}

\begin{fr}
\begin{ansblock}
Jej błąd maleje szybciej, gdy skracasz odległość \dash{} jak $h^{3}$, gdzie
błąd prostej maleje jak $h^{2}$. \emph{Rząd}, a nie rozmiar.
\end{ansblock}

\textbf{To rozróżnienie jest różnicą między twierdzeniem o jednym kroku a
twierdzeniem o wszystkich.} \enquote{Dokładniejsza} to fakt o tym $h$, które
akurat wypróbowałeś; \enquote{o rząd lepsza} mówi, co dzieje się w miarę
zbliżania się, i da się to sprawdzić.

Skrypt sprawdza to właśnie tak: połowi $h$ osiem razy i stwierdza, że błąd
liniowy dzieli się za każdym razem mniej więcej przez
$2^{\val{p17.order.lin}}$, a błąd kwadratowy mniej więcej przez
$2^{\val{p17.order.quad}}$. Przy $h = \val{p17.taylor.h}$ oba błędy wynoszą
$\val{p17.err.lin}$ i $\val{p17.err.quad}$, ale te dwie liczby są
ilustracjami; wynikiem są stosunki.

\begin{note}
To ta sama dyscyplina, której Program~\ref{prog:P05} użył dla swojego
rozrzutu, a Program~\ref{prog:P15} dla swojego kąta zygzaka:
\textbf{stwierdzaj niezmiennik, nie obserwację.} Dwie liczby błędu byłyby
faktem o jednym rozmiarze kroku i jednej funkcji. Rząd przeżywa oba.
\end{note}
\end{fr}

\begin{fr}
Jednej rzeczy ten model nie robi, a powiedzenie tego teraz oszczędza kłopotu
później.

Model Taylora jest \emph{lokalny}. Zapisz, co cię to kosztuje, w kategoriach
obrazu, który narysował Program~\ref{prog:P15}.
\dotline
\nextframe
\end{fr}

\begin{fr}
\begin{ansblock}
Nie mówi nic o tym, gdzie jest minimum. Opisuje powierzchnię w pobliżu punktu
i nie ma zdania o niczym poza nim.
\end{ansblock}

Co jest dokładnie odkryciem Programu~\ref{prog:P15} w innym przebraniu:
gradient jest stwierdzeniem lokalnym, a położenie minimum globalnym. Hesjan
też jest lokalny i dołożenie go tego nie zmienia \dash{} czyni opis lokalny
lepszym, a nie szerszym.

\begin{rigourbox}
Stwierdzenie, że błąd jest $O(h^{3})$, wymaga, by trzecia pochodna istniała i
była ograniczona w pobliżu punktu, co jest twierdzeniem Taylora z resztą i
znajduje się w każdym pierwszym kursie analizy. Ten program używa wniosku i
mierzy go; nie dowodzi go.

Bez dowodu warto wiedzieć, gdzie to zawodzi: w punkcie, w którym funkcja nie
jest gładka, nie ma modelu kwadratowego i cały argument tego programu nie ma
się do czego przyczepić. Program~\ref{prog:P16} podaje ten przypadek: $\relu$ dokładnie w zerze.
\end{rigourbox}
\end{fr}

\begin{fr}
Przed macierzą jedno pytanie, którego odpowiedź już znasz.

W punkcie stacjonarnym gradient jest zerem, więc składnik liniowy znika i
model to sama wartość plus składnik z krzywizną. Co mówi ci tam znak
krzywizny?
\nextframe
\end{fr}

\begin{fr}
\begin{ansblock}
Dodatni to minimum, ujemny maksimum \dash{} i w jednym wymiarze to cała
klasyfikacja.
\end{ansblock}

Program~\ref{prog:F11} powiedział, że zerowa pochodna to nie minimum, i
zostawił test \enquote{pochodnej pochodnej}. To jest ten test i w jednym
wymiarze jest zupełny.

W więcej niż jednym nie jest, a Program~\ref{prog:P10} pokazał już dlaczego:
istnieją kształty bez jednowymiarowego odpowiednika. Przejście od znaku do
klasyfikacji jest tym, do czego służy macierz z następnej sekcji.
\end{fr}

% ============================================================
\section{Macierz drugich pochodnych}
\label{sec:P17-hessian}
\index{hesjan}

\begin{fr}
Przy kilku wejściach jest druga pochodna dla każdej ich \emph{pary}: różniczkuj
po $x_i$, a potem po $x_j$.

Ułóż je w macierz, a jest to \textbf{hesjan}:
\[ H_{ij} = \frac{\partial^{2} f}{\partial x_i \, \partial x_j} \]

Program~\ref{prog:P16} powiedział, że pochodna funkcji wektorowej jest
macierzą z jednym wierszem na wyjście. Powiedz, jaki kształt ma ta i dlaczego
właśnie taki.
\dotline
\nextframe
\end{fr}

\begin{fr}
\begin{ansblock}
Kwadratowy, jeden wiersz i jedna kolumna na wejście \dash{} bo jest to
jakobian \emph{gradientu}, który jest funkcją z wejść w tyle samo wyjść.
\end{ansblock}

Więc i tutaj nie zdefiniowano niczego nowego. Gradient przeprowadza $n$ wejść
w $n$ wyjść, a reguła Programu~\ref{prog:P16} daje jego pochodnej $n$ wierszy
i $n$ kolumn.

Model drugiego rzędu w wielu wymiarach jest wtedy tym jednowymiarowym, w
którym dwa iloczyny zastąpiono dwiema rzeczami, jakich dostarczają
Program~\ref{prog:P05} i Program~\ref{prog:P10}:
\[ f(\vect{p} + \vect{h}) \approx f(\vect{p})
   + \nabla f \cdot \vect{h}
   + \tfrac{1}{2}\, \vect{h}^{\mathsf{T}} H \vect{h} \]
iloczynem skalarnym i formą kwadratową.
\end{fr}

\begin{fr}
Ta forma kwadratowa jest powodem, dla którego ten program może w ogóle mówić o
kierunkach, i opiera się na jednej własności $H$. Której?
\nextframe
\end{fr}

\begin{fr}
\begin{ansblock}
Jest \textbf{symetryczna}: różniczkowanie w obu kolejnościach daje tę samą
odpowiedź.
\end{ansblock}

A symetria jest tym, co odblokowuje wszystko, co udowodnił
Program~\ref{prog:P10}. Macierz symetryczna ma pełny ortogonalny zbiór
wektorów własnych i rzeczywiste wartości własne, więc powierzchnia w pobliżu
punktu ma zbiór prostopadłych kierunków, każdy z własną krzywizną, a
klasyfikacja misa \dash{} siodło \dash{} grzbiet stosuje się bez zmian.

Skrypt sprawdza, że obie pochodne mieszane zgadzają się lepiej niż
$\val{p17.sym.bound}$ w punkcie funkcji, w którą nie wbudowano żadnej symetrii.

\begin{note}
To sprawdzenie jest potwierdzeniem, a nie powodem. Symetria jest twierdzeniem
\dash{} wymaga ciągłości drugich pochodnych cząstkowych w pobliżu punktu i
znajduje się w każdym pierwszym kursie analizy \dash{} a numeryczna zgodność
w jednym punkcie byłaby kiepskim dowodem twierdzenia o każdej funkcji.

Zapisany jest sufit, a nie pomiar, z powodu, za który zapłacił
Program~\ref{prog:P06}: różnica jest szumem zaokrągleń, którego rozmiar jest
własnością maszyny, a tutaj wychodzi dokładnie zero, co jest szczęściem tej
funkcji, a nie czymś, co inna kompilacja jest ci winna.
\end{note}
\end{fr}

\begin{fr}
A więc klasyfikacja, którą Program~\ref{prog:P10} dał formom kwadratowym, jest
teraz stwierdzeniem o \emph{dowolnej} gładkiej funkcji w punkcie stacjonarnym.
Powiedz je: co mówią ci znaki wartości własnych hesjanu?
\dotline
\nextframe
\end{fr}

\begin{fr}
\begin{ansblock}
Same dodatnie to minimum, same ujemne maksimum, znaki mieszane siodło, a
zerowa wartość własna znaczy, że model drugiego rzędu nie rozstrzyga.
\end{ansblock}

Co jest misą Programu~\ref{prog:P10}, jego odwróconą misą, jego siodłem i jego
grzbietem, w tej kolejności, i żadna nowa praca nie była potrzebna, by tu
dojść \dash{} hesjan jest tą macierzą, o której tamten program już mówił.

Teraz użyj tego na wyjaśnieniu, po które sięgają wszyscy, gdy przebieg
przestaje się poprawiać. Model ma $\val{p17.params}$ parametrów, więc punkt
stacjonarny jego straty ma tyle samo znaków wartości własnych. Co musi być o
nich prawdą, żeby punkt był minimum, i na ile sposobów można to chybić o jeden
znak?
\dotline
\nextframe
\end{fr}

\begin{fr}
\begin{ansblock}
Każdy pojedynczy musi być dodatni. Chybić o jeden znak można na
$\val{p17.params}$ sposobów, po jednym na parametr \dash{} a to tylko
sposoby chybienia o jeden.
\end{ansblock}

\begin{trapbox}
\textbf{\enquote{Moja powierzchnia straty ma minima lokalne i to w nich
zacina się trening.}} To obraz dwuwymiarowy przeniesiony tam, gdzie nie
przeżywa, a mówi dlaczego zliczanie znaków: minimum wymaga zgody każdego
znaku, a siodło jednego odszczepieńca, więc punkt stacjonarny w wysokim wymiarze
przytłaczająco prawdopodobnie w ogóle nie jest minimum.

To, w czym zwykle siedzi przebieg, który przestał się poprawiać, jest jedną z
dwóch innych rzeczy. Siodłem, które ma kierunek w dół i gradient zbyt mały, by
za nim szybko pójść. Albo długą wąską doliną, która ma kierunek w dół i jest
przecinana, a nie śledzona \dash{} a to drugie jest tematem następnej sekcji,
a nie żadnym rodzajem minimum.

Uczciwa granica tego argumentu: jest to arytmetyka znaków, a nie pomiar
jakiejkolwiek wytrenowanej sieci, i tego drugiego ta książka nie mierzyła. Co
ustala, to że zerowy gradient jest znacznie słabszym świadectwem minimum w
miliardzie wymiarów niż w dwóch.
\end{trapbox}
\end{fr}

% ============================================================
\section{Dlaczego rozmiar kroku ma sufit}
\label{sec:P17-stepsize}
\index{krzywizna!a rozmiar kroku}

\begin{fr}
Teraz zapłata, a Program~\ref{prog:P15} zostawił ją dokładnie o jedną linijkę
stąd.

Na funkcji kwadratowej spadek gradientu mnoży współrzędną każdego kierunku
własnego przez $1 - \eta\lambda$ w każdym kroku, z własną $\lambda$ tego
kierunku. P15 zmierzył $\val{p15.fac.hi}$ i
$\val{p15.fac.lo}$ na swojej misie i powiedział wprost, że to, gdzie leży
granica, należy do tego programu.

Więc ją znajdź. Co musi spełniać $1 - \eta\lambda$, żeby współrzędna nie
rosła?
\dotline
\nextframe
\end{fr}

\begin{fr}
\begin{ansblock}
$\lvert 1 - \eta\lambda \rvert < 1$, co dla dodatnich $\eta$ i $\lambda$ daje
$\eta < \dfrac{2}{\lambda}$.
\end{ansblock}

\textbf{I to jest całe wyprowadzenie.} Wielokrotne mnożenie przez liczbę
rośnie nieograniczenie dokładnie wtedy, gdy ta liczba przekracza jeden co do
wielkości; liczbą jest tu $1 - \eta\lambda$; a przekształcenie daje rozmiar
kroku.

W tym argumencie nie ma nic o sieciach neuronowych, stratach ani trenowaniu.
To stwierdzenie o ciągu geometrycznym, czyli o tym, co należy do
Programu~\ref{prog:F04}.

\mermaidfig{p17-why-the-step-caps}{Jeden czynnik na kierunek i to, co go
ogranicza.}{sufit na krok}
\end{fr}

\begin{fr}
Teraz część, która czyni z tego stwierdzenie o całej powierzchni, a nie o
jednym kierunku.

Misa ma naraz kilka wartości własnych, a spacer musi pozostać ograniczony w
\emph{każdym} kierunku. Na misie Programu~\ref{prog:P10} obie wynoszą
$\val{p17.lam.hi}$ i $\val{p17.lam.lo}$. Jakie jest największe $\eta$, które
działa, i który kierunek o tym zdecydował?
\dotline
\nextframe
\end{fr}

\begin{fr}
\ans{$\val{p17.eta.star}$, ustalone przez najbardziej stromy}
Bo każdy kierunek żąda $\eta < 2/\lambda$, a wygrywa najmniejsze z tych żądań.
Stromy kierunek pozwala na $\val{p17.eta.star}$; płytki pozwoliłby na
$\val{p17.eta.shallow}$, co jest $\val{p17.eta.gap}$ razy większe i jest bez
znaczenia, bo wzięcie tego wysadziłoby tamten drugi.

\textbf{A więc kierunek o najszybszej krzywiźnie ustala rozmiar kroku dla
całego modelu, a każdy inny kierunek żyje z tym, co dostanie.} To jest zdanie,
którego obrazem był zygzak z Programu~\ref{prog:P15}.

Skrypt nie tylko to wyprowadza. Przemiata $\eta$ i znajduje wartość, przy
której spacer przestaje być zbieżny, a wychodzi ona
$\val{p17.eta.measured}$ \dash{} wyprowadzona granica, zmierzona.
\end{fr}

\begin{fr}
I brama, bo dwa programy opisujące jeden spacer nie mogą się co do niego
różnić. Program~\ref{prog:P15} chodzi po tej misie przy
$\eta = \val{p17.p15.eta}$. Wobec granicy $\val{p17.eta.star}$ \dash{} czy ten
spacer jest bezpieczny?
\nextframe
\end{fr}

\begin{fr}
\begin{ansblock}
Tak, z zapasem. Jego czynnik to
$1 - \val{p17.p15.eta} \times \val{p17.lam.hi} = \val{p17.p15.factor}$:
ujemny, więc stroma współrzędna zmienia stronę w każdym kroku, i mniejszy od
jedynki co do rozmiaru, więc się kurczy.
\end{ansblock}

\textbf{Zmiana znaku nie jest ostrzeżeniem.} Zaczyna się, gdy tylko $\eta$
przekroczy $1/\lambda$, czyli tutaj połowę granicy, więc naprzemienność jest
normalnym zachowaniem na całej górnej połowie zakresu stabilności. Strata,
która oscyluje, nie jest przez to o włos od rozbieżności; może być na połowie
pułapu i całkiem bezpieczna, a rozróżnienie wymaga czynnika, a nie kształtu
krzywej.

Skrypt czyta zapisany rozmiar kroku z P15 i stwierdza, że
leży on ściśle poniżej granicy z tego programu, więc oba nie mogą się
rozjechać.

\begin{trapbox}
\textbf{\enquote{Krok uczenia był za duży} mówi się zwykle jako ocenę, a jest
to nierówność.} Znaczy ona dokładnie $\eta > 2/\lambda_{\max}$, a wielkość po
prawej jest własnością powierzchni straty w punkcie, w którym stoisz, a nie
optymalizatora czy danych.

Przekrocz granicę o krok, a nie jest to łagodne pogorszenie. Oto ta sama misa
przy $\eta = \val{p17.demo.eta}$, o włos powyżej $\val{p17.eta.star}$, przez
$\val{p17.demo.steps}$ kroków:

\transcript{p17-too-big}

Każdy krok mnoży przez liczbę większą co do wielkości od jedynki, więc wzrost
jest geometryczny, a znak się zmienia. \emph{Strata poszła w NaN} to ten ciąg
biegnący dalej, aż skończy się format, co wycenił Program~\ref{prog:P01}.
\end{trapbox}
\end{fr}

\begin{fr}
Z tych samych dwóch liczb wypada jeszcze jedna wielkość, a jest to taka, którą
ta książka spotkała już dwa razy.

Stosunek największej krzywizny do najmniejszej wynosi tu $\val{p17.kappa}$.
Jak nazywa się ta liczba i gdzie się pojawiała?
\dotline
\nextframe
\end{fr}

\begin{fr}
\begin{ansblock}
\textbf{Wskaźnik uwarunkowania}, hesjanu \dash{} ta sama wielkość, którą
Program~\ref{prog:P11} zdefiniował dla macierzy, a Program~\ref{prog:P09}
uważał, by nie mylić jej z wyznacznikiem.
\end{ansblock}

Tam mierzył, jak bardzo rozwiązanie układu wzmacnia błąd. Tutaj mierzy, jak
bardzo rozmiar kroku, na który ci pozwolono, różni się od rozmiaru kroku,
którego chce wolny kierunek, a oba odczyty są tą samą arytmetyką.

Jego konsekwencją jest tempo. Najlepszym pojedynczym rozmiarem kroku dla
funkcji kwadratowej jest $2/(\lambda_{\max} + \lambda_{\min})$ \dash{} zadanie
dodatkowe 2 pyta dlaczego, a krótka odpowiedź brzmi: tam oba skrajne kierunki
są tak samo źle obsłużone. Nawet przy tym kroku, tutaj $\val{p17.best.eta}$,
wolny kierunek kurczy się o czynnik $\val{p17.rate}$ na
krok, więc zamknięcie $99$ procent pozostałej luki zajmuje
$\val{p17.steps.99}$ kroków \dash{} na misie z tylko dwoma kierunkami i przy
dokładnej arytmetyce.

\textbf{To jest koszt wąwozu wypowiedziany jako liczba}, a rośnie on ze
wskaźnikiem uwarunkowania, a nie z rozmiarem zadania.
\end{fr}

% ============================================================
\section{Użycie krzywizny i jego koszt}
\label{sec:P17-newton}
\index{hesjan!metoda Newtona}

\begin{fr}
Jeśli kłopotem jest to, że jeden rozmiar kroku musi obsłużyć kierunki o różnych
krzywiznach, jest oczywista poprawka: dać każdemu kierunkowi własny.

To metoda Newtona \dash{} zamiast kroku przeciw gradientowi rozwiąż
$H\,\vect{s} = -\nabla f$ i zrób krok o $\vect{s}$. Ile kroków zajmuje to
na funkcji kwadratowej?
\nextframe
\end{fr}

\begin{fr}
\ans{$\val{p17.newton.steps}$}
Dokładnie jeden i ląduje w minimum. Dzielenie każdego kierunku przez jego
własną krzywiznę jest właśnie tą poprawką, której zdaniem poprzedniej sekcji
brakowało, a na funkcji kwadratowej model jest funkcją, więc minimum modelu
jest tym prawdziwym.

\textbf{A więc metoda, która rozwiązuje cały problem sekcji 3, istnieje i jest
stroną matematyki z pierwszego roku.} Pytaniem jest tylko, ile kosztuje, a
odpowiedź jest powodem, dla którego nigdy jej nie użyłeś.
\end{fr}

\begin{fr}
Policz. Model ma $\val{p17.params}$ parametrów. Ile wpisów ma jego hesjan i
ile byłoby to w bajtach przy dwóch bajtach na wpis?
\dotline
\nextframe
\end{fr}

\begin{fr}
\begin{ansblock}
$\val{p17.hess.entries}$ wpisów, czyli około $\val{p17.hess.bytes}$ bajtów.
\end{ansblock}

Postaw to obok drabiny rzędów wielkości z Programu~\ref{prog:F01}, a
przekracza każdy jej szczebel poza najwyższym. Jest większe niż liczba tokenów
w korpusie treningowym i większe niż arytmetyka, którą urządzenie robi w
sekundę; jedyne, co jest tam wyżej, to cała arytmetyka jednego wytrenowania
modelu, jakieś sto razy większa. A to porównanie nie jest jak za jak: tamte
wielkości się \emph{wydaje}, przez miesiące i na klastrze, a tu chodzi o
bajty, które trzeba trzymać naraz. Żadna istniejąca pamięć nie jest w zasięgu
wielu rzędów od tego.

Liczba parametrów jest liczbą, którą da się utrzymać; liczba \emph{par}
parametrów nie jest, a robi to podnoszenie do kwadratu.

I przechowanie jej jest tą tanią częścią. Rozwiązanie układu z nią jest
sześcienne, około $\val{p17.hess.solve}$ operacji, co Program~\ref{prog:P09}
wycenił we własnych kategoriach: eliminacja buduje $n$ liczb tam, gdzie odwrotność
buduje $n^{2}$, a tutaj nawet $n$ jest za duże.

\textbf{A więc hesjan jest dokładnie jak jakobian z
Programu~\ref{prog:P16}: nikt go nie tworzy.} Różnica jest taka, że dla
jakobianu iloczyny, których ludzie faktycznie chcą, są tanie, a dla hesjanu to,
czego ludzie chcą \dash{} jego odwrotność zastosowana do gradientu \dash{}
tanie nie jest.
\end{fr}

\begin{fr}
\begin{note}
To nie koniec tematu i błędem byłoby zostawić wrażenie, że informacja drugiego
rzędu jest nie do użycia. Istnieje rodzina metod, które nigdy nie tworzą $H$ i
budują przybliżenie jego działania z gradientów, które już widziały, oraz
metody używające samej przekątnej.

Nazwanie, gdzie przebiega linia, jest bardziej użyteczne niż lista:
\textbf{cokolwiek potrzebuje pamięci $n^{2}$, odpada przy tej skali, a
cokolwiek potrzebuje $n$, wchodzi w grę.} Optymalizatory
Programu~\ref{prog:P20} mieszczą się wszystkie po właściwej stronie tej linii,
a skalowanie na parametr, które kilka z nich robi, jest pomysłem przekątnej
pod inną nazwą.
\end{note}
\end{fr}

\begin{fr}
Jest druga obiekcja wobec metody Newtona, która nie ma nic wspólnego z kosztem,
a znaczy więcej, niż ludzie się spodziewają. Spójrz jeszcze raz na to, co robi
krok: dzieli przez krzywiznę w każdym kierunku.

Co dzieje się w siodle?
\dotline
\nextframe
\end{fr}

\begin{fr}
\begin{ansblock}
Niektóre krzywizny są ujemne, więc dzielenie przez nie robi w tych kierunkach
krok \emph{w stronę} siodła, a nie od niego.
\end{ansblock}

Metoda Newtona znajduje punkty stacjonarne, a nie minima, i nie odróżnia ich
\dash{} jest to wielowymiarowa postać faktu, na który
Program~\ref{prog:F11} wciąż wskazywał, że zerowa pochodna to nie minimum.

\textbf{W połączeniu z argumentem zliczającym z sekcji 2} jest to prawdziwy
problem, a nie ciekawostka: w wysokim wymiarze większość punktów stacjonarnych
to siodła, więc metoda zbieżna do punktów stacjonarnych jest zbieżna do
niewłaściwego rodzaju rzeczy. Każda praktyczna metoda drugiego rzędu musi coś
z tym zrobić, a to, co robi, jest modyfikacją, a nie samą metodą.
\end{fr}

% ============================================================
\section{Ostre i płaskie, uczciwie}
\label{sec:P17-sharpness}
\index{krzywizna!ostre i płaskie minima}

\begin{fr}
Program~\ref{prog:P10} użył słów \emph{wąwóz} i \emph{ostre minimum} dla
niecek o dużym stosunku wartości własnych. Na tych słowach zbudowane jest
znane twierdzenie: że płaskie minima uogólniają się lepiej niż ostre.

Ta sekcja jest o tym, które części tego są stwierdzeniami o wytrenowanym
modelu. Zacznij od arytmetyki. Jeśli przeskalujesz jeden parametr \dash{}
napisz $w = c\,u$ i podziel wejście tego parametru przez $c$ \dash{} co
stanie się z funkcją, którą liczy sieć?
\dotline
\nextframe
\end{fr}

\begin{fr}
\begin{ansblock}
Nic. Każde wyjście jest identyczne dla każdego wejścia, bo dwa czynniki $c$
się skracają.
\end{ansblock}

Model jest tym samym modelem. Jest zapisany w innych współrzędnych, co jest
zamianą bazy z Programu~\ref{prog:P09} zastosowaną do parametru, a nie do
wektora.

Teraz strata jako funkcja $u$, a nie $w$. Jej druga pochodna łapie czynnik
$c^{2}$ z reguły łańcuchowej. Więc przy $c = \val{p17.rescale.c}$ krzywizna
zmienia się o jaki czynnik?
\nextframe
\end{fr}

\begin{fr}
\ans{$\val{p17.rescale.factor}$}
Skrypt sprawdza obie połowy: że wyjścia zgadzają się co do ostatniego bitu przy
każdym wypróbowanym wejściu i że krzywizna jest mnożona dokładnie przez
$\val{p17.rescale.factor}$.

\textbf{A więc jeden model ma dwie ostrości i nic ich nie odróżnia poza tym,
jak ktoś postanowił zapisać parametry.} Liczba, która zmienia się, gdy funkcja
się nie zmienia, jest faktem o układzie współrzędnych.

\mermaidfig{p17-sharp-is-not-a-fact}{Ta sama funkcja, dwie krzywizny.}{ostrość
to nie fakt}
\end{fr}

\begin{fr}
\begin{trapbox}
\textbf{\enquote{To minimum jest płaskie, więc uogólni się lepiej}} nie jest
twierdzeniem, które da się ocenić, dopóki ktoś nie powie, w czym mierzy tę
\emph{płaskość}.

Goła wartość własna hesjanu nią nie jest, na mocy ramki powyżej. Ani żadna
wielkość zbudowana z samych gołych wartości własnych \dash{} ślad, największa,
stosunek \dash{} bo wszystkie one ruszają się przy przeskalowaniu, które nie
zmienia w modelu niczego.

To nie jest techniczny szczegół dotyczący nietypowych architektur. Sieć z
dowolnym skalowaniem po współrzędnych \dash{} czyli prawie każda \dash{}
dopuszcza dokładnie tę reparametryzację, a warstwy normalizujące czynią z niej
symetrię, którą architektura ma z konstrukcji.
\end{trapbox}
\end{fr}

\begin{fr}
Uczciwe stanowisko warto wypowiedzieć w całości, bo łatwo przesadzić w drugą
stronę i orzec, że cały pomysł jest pusty.

To, co tu ustalono, jest wąskie: goła liczba opisująca ostrość nie jest
własnością funkcji. To nie obala twierdzenia empirycznego i nie znaczy, że nic
nie przeżywa.

Czego trzeba, by twierdzenie oparte na ostrości było o modelu, a nie o układzie
współrzędnych?
\dotline
\nextframe
\end{fr}

\begin{fr}
\begin{ansblock}
Miary, która przekształca się tak samo jak parametry, tak by przeskalowanie
zostawiało ją bez zmian.
\end{ansblock}

Co jest rzeczą, o którą naprawdę można prosić, i rzeczą, którą naprawdę da się
sprawdzić, a jest to pytanie, na które staranna praca na ten temat odpowiada
wprost, a niestaranna go nie wspomina.

\begin{warning}
\textbf{Ta książka nie zmierzyła twierdzenia o uogólnianiu i go nie
rozstrzyga.} Zrobienie tego wymaga wytrenowanych modeli, definicji płaskości,
która przeżywa ramkę powyżej, i dość przebiegów, by pokonać wariancję
\dash{} co jest tematem Programu~\ref{prog:P34}, a nie tego.

Co ten program ustala, to filtr, a nie werdykt: mając twierdzenie o ostrości,
zapytaj, czy wielkość, którą ono nazywa, zmieniłaby się przy przeskalowaniu
zostawiającym model identycznym. Jeśli tak, twierdzenie jest o układzie
współrzędnych. To pytanie, na które odpowiesz w minutę, i załatwia ono dużą
część tego, co się mówi.
\end{warning}
\end{fr}

\begin{fr}
Jeszcze jedna obserwacja i domyka ona pętlę z powrotem do sekcji 3.

Ograniczenie na krok $\eta < 2/\lambda_{\max}$ ma ten sam problem, a zauważenie
tego jest uczciwym sposobem utrzymania obu rzeczy naraz.

Powiedz dlaczego, a potem powiedz, dlaczego dla tamtego użycia to nie ma
znaczenia.
\dotline
\nextframe
\end{fr}

\begin{fr}
\begin{ansblock}
Bo $\lambda_{\max}$ też się przeskalowuje. Przy przeskalowaniu
\emph{jednorodnym} nie ma to znaczenia, bo $\eta$ przeskalowuje się razem z nią
w przeciwną stronę i \emph{krok faktycznie zrobiony} pozostaje bez zmian.
\end{ansblock}

\textbf{Jednorodne wykonuje tu prawdziwą pracę.} Przeskaluj jeden parametr
spośród wielu, czyli przypadek zbudowany w sekcji 5, a jeden wspólny $\eta$ nie
da rady za nim pójść: $\lambda_{\max}$ się przesuwa, a krok każdego innego
kierunku przesuwa się razem z nim. To przypadek z zadania dodatkowego 4
przeciw temu tutaj, a różnica między nimi polega na tym, czy zmiana układu
współrzędnych jest zmianą \emph{każdej} współrzędnej.

To jest różnica między tymi dwoma użyciami w jednym zdaniu. W sekcji 3
krzywizna pojawia się obok rozmiaru kroku i para jest niezmiennicza, choć żaden
z jej członków nie jest. W sekcji 5 krzywizna jest cytowana sama, a sama nie
znaczy nic.

\textbf{A więc właściwym pytaniem o każdą liczbę opisującą krzywiznę nigdy nie
jest \enquote{czy jest duża}, tylko \enquote{względem czego jest duża}}, a
sekcja 3 ma na to odpowiedź tam, gdzie literatura o ostrości musi ją dopiero
dostarczyć.
\end{fr}

% ============================================================
\begin{summarybox}
\sumitem{1--2}{\result{\textbf{Model Taylora drugiego rzędu to wartość, nachylenie i
  połowa krzywizny razy kwadrat}, a $\tfrac{1}{2}$ jest tym, co sprawia, że
  druga pochodna modelu równa się drugiej pochodnej funkcji}.}
\sumitem{3--4}{\result{\textbf{\enquote{Najlepszy} znaczy rząd, a nie rozmiar}: błąd
  kwadratowego maleje jak $h^{\val{p17.order.quad}}$, gdzie błąd prostej maleje
  jak $h^{\val{p17.order.lin}}$, co sprawdza się połowieniem $h$, a nie
  cytowaniem dwóch błędów}.}
\sumitem{5--6}{\result{\textbf{Model jest lokalny.} Czyni opis lokalny lepszym, a nie
  szerszym, i nie mówi nic o tym, gdzie jest minimum}.}
\sumitem{7--8}{\result{W punkcie stacjonarnym składnik liniowy znika, więc znak
  krzywizny klasyfikuje \dash{} zupełnie w jednym wymiarze, a w większej
  liczbie nie}.}
\sumitem{9--10}{\result{\textbf{Hesjan jest jakobianem gradientu}, więc reguła
  Programu~\ref{prog:P16} czyni go kwadratowym, z jednym wierszem i jedną
  kolumną na wejście}.}
\sumitem{11--12}{\result{\textbf{Jest symetryczny}, więc Program~\ref{prog:P10} stosuje
  się bez zmian: rzeczywiste wartości własne, prostopadłe kierunki, każdy z
  własną krzywizną}.}
\sumitem{13--14}{\result{Znaki wartości własnych dają misę, odwróconą misę, siodło i
  grzbiet \dash{} klasyfikację z P10, teraz dla dowolnej gładkiej funkcji w
  punkcie stacjonarnym}.}
\sumitem{15}{\result{\textbf{Punkt stacjonarny w wysokim wymiarze przytłaczająco
  prawdopodobnie nie jest minimum}: minimum wymaga, by każdy z
  $\val{p17.params}$ znaków był dodatni, a chybić o jeden znak można na tyle
  samo sposobów. \enquote{Zacięcie w minimum lokalnym} to obraz dwuwymiarowy
  przeniesiony tam, gdzie nie przeżywa}.}
\sumitem{16--17}{\result{\textbf{Spacer pozostaje ograniczony dokładnie wtedy, gdy
  $\lvert 1 - \eta\lambda\rvert < 1$, czyli $\eta < 2/\lambda$}, co jest faktem
  o ciągu geometrycznym i o niczym więcej}.}
\sumitem{18--19}{\result{\textbf{Najbardziej stromy kierunek ustala krok dla całego
  modelu.} Tutaj $\val{p17.eta.star}$ wobec $\val{p17.eta.shallow}$ kierunku
  płytkiego, czyli czynnik $\val{p17.eta.gap}$. Wyprowadzone, a potem zmierzone
  przemiataniem: $\val{p17.eta.measured}$}.}
\sumitem{20--21}{\result{Spacer z Programu~\ref{prog:P15} przy $\val{p17.p15.eta}$
  mieści się pod granicą, a skrypt to stwierdza. \textbf{\enquote{Krok uczenia
  był za duży} to nierówność $\eta > 2/\lambda_{\max}$}, a za nią wzrost jest
  geometryczny}.}
\sumitem{22--23}{\result{\textbf{Wskaźnik uwarunkowania $\val{p17.kappa}$ to ta sama
  wielkość, którą zdefiniował Program~\ref{prog:P11}}, odczytana na hesjanie.
  Przy najlepszym kroku wciąż potrzeba $\val{p17.steps.99}$ kroków, by
  zamknąć $99$ procent luki wolnego kierunku}.}
\sumitem{24--25}{\result{\textbf{Metoda Newtona ląduje w minimum funkcji kwadratowej w
  $\val{p17.newton.steps}$ kroku}, dając każdemu kierunkowi własny rozmiar
  kroku}.}
\sumitem{26--28}{\result{\textbf{A jej hesjan ma $\val{p17.hess.entries}$ wpisów przy
  realnym rozmiarze modelu}, z sześciennym rozwiązywaniem układu na dokładkę.
  Hesjanu też nikt nie tworzy \dash{} i w odróżnieniu od jakobianu iloczyn,
  którego ludzie chcą, nie jest tani}.}
\sumitem{29--30}{\result{\textbf{Metoda Newtona znajduje punkty stacjonarne, a nie
  minima}, więc w siodle robi krok w złą stronę w kierunkach ujemnych. Wraz z
  argumentem zliczającym z sekcji 2 jest to prawdziwy problem w wysokim
  wymiarze}.}
\sumitem{31--33}{\result{\textbf{Przeskalowanie parametru przez $c$ zmienia
  krzywiznę o $c^{2}$ \dash{} $\val{p17.rescale.factor}$ przy tym $c$,
  którego używają ramki \dash{} i zostawia każde wyjście identyczne.} Więc
  jeden
  model ma dwie ostrości, a goła liczba opisująca ostrość jest faktem o
  układzie współrzędnych}.}
\sumitem{34}{\result{\textbf{Żadna wielkość zbudowana z gołych wartości własnych tego
  nie przeżywa} \dash{} ani największa, ani ślad, ani stosunek \dash{} a
  warstwy normalizujące czynią z przeskalowania symetrię, którą architektura ma
  z konstrukcji}.}
\sumitem{35--36}{\result{Przeżyłaby miara przekształcająca się tak jak parametry.} \textbf{Ta książka nie zmierzyła twierdzenia o uogólnianiu i go nie
  rozstrzyga}; dostarcza filtru.}
\sumitem{37--38}{\result{\textbf{Ograniczenie na krok ma ten sam problem i przy
  przeskalowaniu \emph{jednorodnym} nie cierpi z jego powodu}, bo $\eta$
  przeskalowuje się razem z $\lambda$, a krok zrobiony pozostaje bez zmian;
  przeskaluj jeden parametr spośród wielu, a jeden wspólny $\eta$ nie da rady
  za nim pójść. Pytaniem o krzywiznę nigdy nie jest \emph{czy jest duża},
  tylko \emph{względem czego}}.}
\end{summarybox}

\canyou

\begin{testexercises}
\item Funkcja kwadratowa ma krzywiznę $\lambda = 8$ w swoim jedynym kierunku.
  Podaj największy rozmiar kroku, który jest zbieżny, i powiedz, co dzieje się
  dokładnie przy tej wartości.
  \answerto{$\eta < 2/8 = \num{0.25}$. Dokładnie przy $\num{0.25}$ czynnik
  wynosi $1 - 2 = -1$, więc współrzędna zmienia znak i zachowuje swój rozmiar
  na zawsze \dash{} ani nie jest zbieżna, ani rozbieżna, co jest granicą będącą
  granicą, a nie marginesem bezpieczeństwa.}
\item Strata ma w bieżącym punkcie wartości własne hesjanu $50$, $5$ i
  $\num{0.5}$. Podaj pułap na krok, wskaźnik uwarunkowania i to, który
  kierunek porusza się w efekcie najwolniej.
  \answerto{$\eta < 2/50 = \num{0.04}$, ustalone przez $50$. Wskaźnik
  uwarunkowania to $50/\num{0.5} = 100$. Problemem jest kierunek
  $\num{0.5}$: wolno mu na $\eta < 4$, a dostaje $\num{0.04}$, więc porusza się
  w setnej części tempa, na jakie mógłby.}
\item Dlaczego wyprowadzenie z sekcji 3 nigdy nie wspomina, ile parametrów ma
  model?
  \answerto{Bo jest to stwierdzenie o jednym kierunku własnym naraz, a
  kierunki nie oddziałują na funkcji kwadratowej. Dołożenie kierunków dokłada
  żądania tej samej postaci, a wiążącym wciąż jest największa krzywizna. Liczba
  parametrów zmienia to, jakie $\lambda_{\max}$ jest prawdopodobne; do
  argumentu nie wchodzi.}
\item Ktoś donosi, że jego model zbiegł do minimum o największej wartości
  własnej hesjanu $\num{0.02}$, i wnioskuje, że jest to płaskie minimum, które
  dobrze się uogólni. Podaj dwa pytania, które byś zadał.
  \answerto{Pierwsze: w jakich współrzędnych? Przeskalowanie parametru rusza tę
  liczbę o kwadrat skali, gdy model pozostaje bez zmian, więc samo
  $\num{0.02}$ nie mówi nic. Drugie: względem czego? Krzywizna ma sens tylko
  wobec czegoś o tych samych jednostkach, czym w sekcji 3 jest rozmiar kroku.
  Żadne z pytań nie jest wrogie i oba mają odpowiedzi; twierdzenie, które nie
  potrafi na nie odpowiedzieć, jest o układzie współrzędnych.}
\item Twój przebieg treningowy się rozbiega. Podaj stwierdzenie arytmetyczne
  odpowiadające \enquote{krok uczenia był za duży} i jedną rzecz inną niż
  obniżenie $\eta$, która by to naprawiła.
  \answerto{$\eta > 2/\lambda_{\max}$ w punkcie, w którym przebieg był. Inaczej
  niż obniżając $\eta$: zmniejszyć $\lambda_{\max}$, co robią i normalizacja, i
  mniejsza inicjalizacja, skoro ograniczenie dotyczy powierzchni tak samo jak
  optymalizatora. Zauważ, że $\lambda_{\max}$ jest własnością tego, gdzie
  stoisz, więc przebieg może być stabilny przez tysiąc kroków, a potem nie.}
\end{testexercises}

\begin{furtherproblems}
\item Pokaż, że na funkcji kwadratowej o krzywiźnie $\lambda$ krok
  $\eta = 1/\lambda$ trafia w minimum jednym ruchem.
  \answerto{Czynnik wynosi $1 - \eta\lambda = 0$, więc współrzędna idzie do
  zera w jednym kroku. To metoda Newtona w jednym wymiarze: dzielenie przez
  krzywiznę jest dokładnie wyborem $\eta = 1/\lambda$, a powodem, dla którego
  nie uogólnia się tanio, jest to, że w wielu wymiarach dzielenie jest
  rozwiązywaniem układu.}
\item Dlaczego $\eta = 2/(\lambda_{\max} + \lambda_{\min})$ jest najlepszym
  pojedynczym rozmiarem kroku dla funkcji kwadratowej, a nie coś bliższego
  $1/\lambda_{\max}$?
  \answerto{Bo najgorszym kierunkiem jest ten ze skrajnych, który ma większe
  $\lvert 1 - \eta\lambda\rvert$, a ten najgorszy przypadek jest minimalizowany
  tam, gdzie oba są równe i przeciwne. Rozwiązanie
  $1 - \eta\lambda_{\min} = \eta\lambda_{\max} - 1$ je daje. Wynikające tempo
  to $(\kappa - 1)/(\kappa + 1)$, skąd wzięło się $\val{p17.rate}$.}
\item Hesjan sumy jest sumą hesjanów. Co to mówi o krzywiźnie straty
  uśrednionej po wsadzie?
  \answerto{Że jest średnią krzywizn poszczególnych przykładów, więc
  $\lambda_{\max}$ wsadu jest ograniczone przez największe z pojedynczych i
  zwykle jest znacznie mniejsze, bo poszczególne hesjany nie dzielą swoich
  stromych kierunków. To jeden uczciwy powód, dla którego większy wsad toleruje
  większy krok \dash{} i jest to stwierdzenie o powierzchni, a nie o szumie, co
  należy do Programu~\ref{prog:P21}.}
\item Przeskalowanie z sekcji 5 użyło jednego parametru. Co ten sam argument
  robi z całym hesjanem, jeśli przeskalujesz każdy parametr przez $c$?
  \answerto{Mnoży go przez $c^{2}$ w całości, więc każda wartość własna się
  skaluje, a wskaźnik uwarunkowania \dash{} \emph{stosunek} \dash{} nie. Warto
  to zauważyć: wskaźnik uwarunkowania jest niezmienniczy przy jednolitym
  przeskalowaniu tam, gdzie pojedyncze wartości własne nie są, co jest jednym z
  powodów, dla których jest wielkością bardziej cytowalną. Nie jest
  niezmienniczy przy różnym skalowaniu parametrów, a taki przypadek zbudowała
  sekcja 5.}
\item Dlaczego zerowa wartość własna hesjanu zostawia klasyfikację
  nierozstrzygniętą, zamiast znaczyć, że kierunek jest płaski?
  \answerto{Bo model drugiego rzędu nic tam nie mówi, a zachowanie funkcji w
  tym kierunku rozstrzygają składniki, które model porzucił. $x^{4}$ i $-x^{4}$
  mają w zerze tę samą zerową drugą pochodną, a jedna ma minimum tam, gdzie
  druga ma maksimum. To granica modelu lokalnego pokazująca się sama i dlatego
  \enquote{płaski kierunek} i \enquote{zerowa krzywizna} to nie to samo
  stwierdzenie.}
\end{furtherproblems}
